\documentclass{article}
\usepackage{graphicx}
\usepackage{amsthm}
\usepackage{amssymb}
\usepackage{amsmath}
\usepackage{parskip}
\usepackage[margin=0.6in]{geometry}
\usepackage{tikz}
\usepackage{titling}

 % move title upward

\title{15.095 Machine Learning HW2}
\author{Chris Low}
\date{} % removes the date

\begin{document}

\maketitle

\section*{Question 1}

\subsection*{(a)}
\textbf{False.}

Consider the underlying function
\[
f(x) = \sin x,
\]
which is non-convex on $\mathbb{R}$. Let $n=2$ and choose
\[
x_1 = 0, \qquad x_2 = \frac{\pi}{2}.
\]
Then
\[
y_1 = f(x_1) = 0, \qquad y_2 = f(x_2) = 1.
\]

So we can set: 
\[
\theta_1 = 0, \qquad \theta_2 = 1,
\qquad
\xi_1 = \xi_2 = \frac{1-0}{\frac{\pi}{2}-0} = \frac{2}{\pi}.
\]
The convexity constraints are satisfied since
\[
\theta_1 + \xi_1(x_2-x_1)
= 0 + \frac{2}{\pi}\left(\frac{\pi}{2}\right)
= 1
= \theta_2,
\]
and
\[
\theta_2 + \xi_2(x_1-x_2)
= 1 + \frac{2}{\pi}\left(-\frac{\pi}{2}\right)
= 0
= \theta_1.
\]

Moreover, since $\theta_i = y_i$ for both observations,
\[
\frac{1}{2}\sum_{i=1}^2 (y_i-\theta_i)^2 = 0.
\]

Thus, even though the underlying function $f$ is non-convex, the convex
regression problem only needs to fit the observed points. These two
sampled points lie on the affine function $g(x)=\frac{2}{\pi}x$, which
is convex. Hence, they can be fit perfectly, giving an objective value
of zero. Therefore, the statement is false. \qed

\newpage
\subsection*{(b)}
\textbf{True.}

Let
\[
r_i = y_i - x_i^\top \beta, \qquad i=1,\ldots,n.
\]

Introduce binary variables $z_i \in \{0,1\}$, where $z_i=1$ means
that residual $i$ is selected. We also introduce continuous variables
$u_i \geq 0$. We can formulate the problem as

\[
\min_{\beta,r,u,z} \sum_{i=1}^n u_i
\]

subject to

\[
r_i = y_i - x_i^\top \beta,
\qquad i=1,\ldots,n,
\]

\[
u_i \geq r_i - M(1-z_i),
\qquad i=1,\ldots,n,
\]

\[
u_i \geq -r_i - M(1-z_i),
\qquad i=1,\ldots,n,
\]

\[
\sum_{i=1}^n z_i = k,
\]

\[
u_i \geq 0,
\qquad i=1,\ldots,n,
\]

\[
z_i \in \{0,1\},
\qquad i=1,\ldots,n,
\]

where $M$ is a sufficiently large constant.

If $z_i=1$, the two Big-$M$ constraints reduce to
\[
u_i \geq r_i,
\qquad
u_i \geq -r_i,
\]
which is equivalent to 
\[
u_i \geq |r_i|.
\]

If $z_i=0$, the Big-$M$ constraints become
\[
u_i \geq r_i-M,
\qquad
u_i \geq -r_i-M.
\]
For sufficiently large $M$, both right-hand sides are less than or
equal to zero. Thus, together with $u_i\geq 0$, the smallest feasible
value is
$
u_i=0.$
Since the objective minimizes $\sum_{i=1}^n u_i$, the optimizer will
set $u_i=0$ for these unselected residuals.


Since
$
\sum_{i=1}^n z_i=k,
$
exactly $k$ residuals must be selected. 


For each selected residual,
$u_i \geq |r_i|$, and since we minimize $\sum_{i=1}^n u_i$, we have
$u_i=|r_i|$ at optimality. 

For each unselected residual, $u_i=0$.
Therefore, minimizing $\sum_{i=1}^n u_i$ selects and sums the $k$
smallest absolute residuals. All constraints and the objective are
linear, with $z_i$ binary, so this is a linear MIO formulation. \qed



\newpage
\subsection*{(c)}


\textbf{False.}

The two constraints do different things. For two highly correlated
variables, the constraint
\[
z_i + z_j \leq 1
\]
means that we can select at most one of them. So the possible choices
are
\[
(z_i,z_j) \in \{(0,0),(1,0),(0,1)\}.
\]

On the other hand, if we put $i$ and $j$ in the same group and impose
\[
z_i=z_j,
\]
then they must either both be selected or both be excluded. The possible
choices are therefore
\[
(z_i,z_j) \in \{(0,0),(1,1)\}.
\]

These are not equivalent. In fact, $z_i=z_j$ allows both highly
correlated variables to be selected at the same time, while
$z_i+z_j\leq1$ is specifically used to prevent this. Therefore, we
cannot replace the pairwise-collinearity constraint with the
group-sparsity constraint. \qed




\vspace{5mm}
\subsection*{(d)}
\textbf{False.}

From the robust logistic regression result in Lecture 6, if the feature
uncertainty set is
\[
\mathcal{U}_x
=
\left\{
\Delta X:
\|\Delta x_i\|_q \leq \rho,
\quad i=1,\ldots,n
\right\},
\]
then the robust counterpart is
\[
\max_{\beta,\beta_0}
-\sum_{i=1}^n
\log\left(
1+
e^{-y_i(\beta^\top x_i+\beta_0)
+\rho\|\beta\|_{q^*}}
\right),
\]
where $q^*$ is the dual norm of $q$, satisfying
\[
\frac{1}{q}+\frac{1}{q^*}=1.
\]

Thus, the robustness term is generally
\[
\rho\|\beta\|_{q^*},
\]
not
\[
\rho\|\beta\|_q.
\]

For example, take
\[
q=\infty,
\qquad
\rho=1,
\qquad
\beta=(1,1).
\]
The dual norm of $\ell_\infty$ is $\ell_1$, so the correct robustness
term is
\[
\rho\|\beta\|_1
=
1(1+1)
=
2.
\]
However, the expression in the statement would use
\[
\rho\|\beta\|_\infty
=
1\cdot 1
=
1.
\]

Therefore, the stated robust counterpart underestimates the worst-case
effect and is not correct in general. The two expressions coincide only
for self-dual norms, such as $q=2$. \qed


\newpage

\vspace{5mm}
\subsection*{(e)}

\textbf{True.}

Let
\[
m_i = y_i(\beta^\top x_i+\beta_0).
\]
The robust problem: 
\[
\max_{\beta,\beta_0}
\min_{\Delta y\in\mathcal U_y}
-\sum_{i=1}^n
\log\left(
1+e^{-(1-2\Delta y_i)m_i}
\right),
\]
where
\[
\mathcal U_y
=
\left\{
\Delta y\in\{0,1\}^n:
\sum_{i=1}^n \Delta y_i\le \Gamma
\right\}.
\]

For fixed $\beta,\beta_0$, we define
\[
f_i(\Delta y_i)
=
-\log\left(
1+e^{-(1-2\Delta y_i)m_i}
\right).
\]
Since $\Delta y_i\in\{0,1\}$, we can write
\[
f_i(\Delta y_i)
=
f_i(0)
+
\bigl(f_i(1)-f_i(0)\bigr)\Delta y_i.
\]


\[
f_i(0)
=
-\log(1+e^{-m_i}),
\]

\[
f_i(1)-f_i(0)
=
\log\left(
\frac{1+e^{-m_i}}{1+e^{m_i}}
\right).
\]
Thus the inner problem becomes
\[
\min_{\Delta y}
-\sum_{i=1}^n\log(1+e^{-m_i})
+
\sum_{i=1}^n c_i\Delta y_i,
\]
where
\[
c_i
=
\log\left(
\frac{1+e^{-m_i}}{1+e^{m_i}}
\right),
\]
subject to
\[
\Delta y_i\in\{0,1\},
\qquad
\sum_{i=1}^n\Delta y_i\le\Gamma.
\]

Now we observe that 
$$
0\le\Delta y_i\le1,
\qquad
\sum_{i=1}^n\Delta y_i\le\Gamma
$$
has integer extreme points when $\Gamma$ is integer, we can therefore relax the binary
restriction, which would not change the optimal value. So 
the inner problem is equivalent to the LP
\[
\min_{\Delta y}
\sum_{i=1}^n c_i\Delta y_i
\]
subject to
\[
\sum_{i=1}^n\Delta y_i\le\Gamma,
\]
\[
\Delta y_i\le1,
\qquad i=1,\ldots,n,
\]
\[
\Delta y_i\ge0,
\qquad i=1,\ldots,n.
\]


Now we can take the dual problem. By strong duality, the optimal value is equal to that of its dual problem. 

Let $\mu\le0$ be the dual variable associated with
\[
\sum_{i=1}^n\Delta y_i\le\Gamma,
\]
and let $\nu_i\le0$ be the dual variable associated with
\[
\Delta y_i\le1.
\]
The dual is
\[
\max_{\mu,\nu}
\Gamma\mu+\sum_{i=1}^n\nu_i
\]
subject to
\[
\mu+\nu_i\le c_i,
\qquad i=1,\ldots,n,
\]
\[
\mu\le0,
\qquad
\nu_i\le0.
\]

By strong duality, the robust problem is therefore equivalent to
\[
\max_{\beta,\beta_0,\mu,\nu}
-\sum_{i=1}^n
\log\left(
1+e^{-y_i(\beta^\top x_i+\beta_0)}
\right)
+
\Gamma\mu
+
\sum_{i=1}^n\nu_i
\]
subject to
\[
\mu+\nu_i
\le
\log\left(
\frac{
1+e^{-y_i(\beta^\top x_i+\beta_0)}
}{
1+e^{y_i(\beta^\top x_i+\beta_0)}
}
\right),
\qquad i=1,\ldots,n,
\]
\[
\mu\le0,
\qquad
\nu_i\le0.
\]

Finally,
\[
\log\left(
\frac{1+e^{-m_i}}{1+e^{m_i}}
\right)
=
-m_i,
\]
so the constraints may equivalently be written as
\[
\mu+\nu_i
\le
-y_i(\beta^\top x_i+\beta_0).
\]

The robust counterpart is therefore a single continuous concave
maximization problem, equivalently a convex optimization problem, and
contains no integer variables. Therefore the statement is true. \qed



\newpage






\section*{Question 2}


\vspace{5mm}
\subsection*{(a)}

Fix $(\beta,\beta_0)$, so the losses $c_i$ are constants, and we order them such that
\[
c_{(1)} \geq c_{(2)} \geq \cdots \geq c_{(n)}.
\]

For $z\in\mathcal Z$, exactly $k$ entries of $z$ equal $1$. Therefore,
the objective is maximized by selecting the $k$ largest losses. Hence, it is equivalent to saying
\[
\max_{z\in\mathcal Z}\sum_{i=1}^n z_i c_i
=
\sum_{j=1}^k c_{(j)}.
\]

Now consider the continuous relaxation
\[
\mathcal P
=
\left\{
z\in\mathbb R^n:
\sum_{i=1}^n z_i=k,\;
0\leq z_i\leq 1
\right\}.
\]

Suppose $z\in\mathcal P$ is optimal. However, there exists indices $a,b$ such that
\[
c_a>c_b,
\qquad
z_a<1,
\qquad
z_b>0.
\]
That is, some positive weight is placed on a smaller loss $c_b$ while
the larger loss $c_a$ has not yet received its maximum weight of $1$.

We define 
\[
\epsilon=\min\{1-z_a,z_b\}>0,
\]

essentially satisfying the conditions that $z_a$ cannot exceed 1, and $z_b$ cannot go below 0. 

We then construct a new solution $z'$ by moving $\epsilon$ units of weight from
$b$ to $a$:
\[
z_a'=z_a+\epsilon,
\qquad
z_b'=z_b-\epsilon,
\]
with all other components unchanged.

Since the same amount of weight is added and removed,
\[
\sum_{i=1}^n z_i'
=
\sum_{i=1}^n z_i
=
k.
\]
Also, with our choice of $\epsilon$, 
\[
0\leq z_i'\leq1.
\]
Therefore, $z'\in\mathcal P$.

Our change in the objective is
\[
\sum_{i=1}^n z_i'c_i-\sum_{i=1}^n z_ic_i
=
(z_a+\epsilon)c_a+(z_b-\epsilon)c_b
-z_ac_a-z_bc_b =
\epsilon c_a-\epsilon c_b
=
\epsilon(c_a-c_b).
\]


Since $\epsilon>0$ and $c_a>c_b$,
\[
\epsilon(c_a-c_b)>0.
\]
Thus, moving weight from the smaller loss to the larger loss strictly
increases the objective, contradicting our initial assumption that $z$ was
optimal.

Therefore, at an optimal point, no positive weight can be placed on a smaller
loss while a larger loss has weight less than $1$. So there exists
an optimal solution that assigns weight $1$ to the $k$ largest losses
and weight $0$ to all remaining losses. In other words, 
\[
\max_{z\in\mathcal P}\sum_{i=1}^n z_i c_i
=
\sum_{j=1}^k c_{(j)}.
\]

We combine the two results,
\[
\max_{z\in\mathcal Z}\sum_{i=1}^n z_i c_i
=
\max_{z\in\mathcal P}\sum_{i=1}^n z_i c_i
=
\sum_{j=1}^k c_{(j)}.
\] \qed


\vspace{5mm}
\subsection*{(b)}


By part (a), for fixed $(\beta,\beta_0)$, we can replace the inner
maximization over $\mathcal Z$ by its continuous relaxation $\mathcal P$.
Let
\[
c_i=\ell(y_i,\beta^\top x_i+\beta_0).
\]
Then the inner problem is
\[
\begin{aligned}
\max_{z}\quad
& \sum_{i=1}^n c_i z_i\\
\text{s.t.}\quad
& \sum_{i=1}^n z_i=k,\\
& z_i\leq 1, \qquad i=1,\ldots,n,\\
& z_i\geq 0, \qquad i=1,\ldots,n.
\end{aligned}
\]

Associate the free dual variable $\theta$ with the  constraint (due to the equality, the dual variable is free)
\[
\sum_{i=1}^n z_i=k,
\]
and the dual variable $u_i\geq0$ with each constraint $z_i\leq1$.

Since each primal variable $z_i\geq0$, the corresponding dual constraint is
\[
\theta+u_i\geq c_i,
\qquad i=1,\ldots,n.
\]

Therefore, the dual of the inner problem is
\[
\begin{aligned}
\min_{\theta,u}\quad
& k\theta+\sum_{i=1}^n u_i\\
\text{s.t.}\quad
& \theta+u_i\geq c_i,
\qquad i=1,\ldots,n,\\
& u_i\geq0,
\qquad i=1,\ldots,n,\\
& \theta\in\mathbb R.
\end{aligned}
\]

By strong duality, the inner maximization and its dual have the same
optimal value. Substituting
\[
c_i=\ell(y_i,\beta^\top x_i+\beta_0)
\]
and combining the dual with the outer minimization over
$(\beta,\beta_0)$ gives
\[
\boxed{
\begin{aligned}
\min_{\beta,\beta_0,\theta,u}\quad
& k\theta+\sum_{i=1}^n u_i\\
\text{s.t.}\quad
& \theta+u_i
\geq
\ell(y_i,\beta^\top x_i+\beta_0),
\qquad i=1,\ldots,n,\\
& u_i\geq0,
\qquad i=1,\ldots,n,\\
& \theta\in\mathbb R.
\end{aligned}}
\]

Thus, problem (2.1) is equivalent to formulation (2.2). \qed

\newpage

Now let the loss be the misclassification loss $
\mathbf{1}\left\{
y_i(\beta^\top x_i+\beta_0)\leq 0
\right\}.
$

Here, the indicator function equals $1$ when its condition is true and
$0$ otherwise. Thus,
\[
y_i(\beta^\top x_i+\beta_0)\leq 0
\]
means that observation $i$ is misclassified, so its loss is $1$.
If
\[
y_i(\beta^\top x_i+\beta_0)>0,
\]
the observation is correctly classified and its loss is $0$.

Since we first define $c_i=\ell(y_i,\beta^\top x_i+\beta_0)$ and the duality argument holds for any fixed values of $c_i$, it remains valid when $\ell$ is the misclassification loss. Therefore, \textbf{formulation (2.2) is still valid} as 
\[
c_i
=
\mathbf{1}\left\{
y_i(\beta^\top x_i+\beta_0)\leq 0
\right\}.
\]

However, \textbf{(2.2) \textbf{cannot be handed directly}} to an LP or convex optimization
solver. In particular, its constraints become
\[
\theta+u_i
\geq
\mathbf{1}\left\{
y_i(\beta^\top x_i+\beta_0)\leq 0
\right\}.
\]
The indicator function jumps from $0$ to $1$ when
$y_i(\beta^\top x_i+\beta_0)$ crosses zero, so it is a discontinuous and
nonconvex function of $(\beta,\beta_0)$. Therefore, the resulting problem
is not directly an LP or convex optimization problem. We might need some further
mixed-integer reformulation, for example we can use binary variables and
big-$M$ constraints, before we give it to a solver.


\vspace{5mm}

\subsection*{(c)}


For the hinge loss,
\[
\ell(y_i,\beta^\top x_i+\beta_0)
=
\max\left\{
1-y_i(\beta^\top x_i+\beta_0),\,0
\right\}.
\]

Substituting this into (2.2), we have
\[
\theta+u_i
\geq
\max\left\{
1-y_i(\beta^\top x_i+\beta_0),\,0
\right\}.
\]

 We see how a constraint in the form
$
a\geq\max\{b,0\}
$
is equivalent to
\[
a\geq b
\qquad\text{and}\qquad
a\geq0.
\]

Therefore, the hinge-loss constraint can be written as
\[
\theta+u_i
\geq
1-y_i(\beta^\top x_i+\beta_0)
\]
and
\[
\theta+u_i\geq0.
\]

Rearrange the first constraint: 
\[
\theta+u_i+y_i x_i^\top\beta+y_i\beta_0\geq1.
\]

Thus, (2.2) becomes
\[
\boxed{
\begin{aligned}
\min_{\beta,\beta_0,\theta,u}\quad
& k\theta+\sum_{i=1}^n u_i\\
\text{s.t.}\quad
& \theta+u_i+y_i x_i^\top\beta+y_i\beta_0\geq1,
\qquad i=1,\ldots,n,\\
& \theta+u_i\geq0,
\qquad i=1,\ldots,n,\\
& u_i\geq0,
\qquad i=1,\ldots,n.
\end{aligned}}
\]

The objective and all constraints are linear, so this is a linear
optimization problem (LP).


\newpage
There are $p$ variables in $\beta$, one variable $\beta_0$, one variable
$\theta$, and $n$ variables $u_i$. Therefore, there are
\[
\boxed{p+n+2}
\]
continuous variables.

There are $n$ constraints of the first type, $n$ constraints of the
second type, and $n$ nonnegativity constraints on $u_i$, giving
\[
\boxed{3n}
\]
constraints in total. Equivalently, there are $2n$ general linear
constraints and $n$ nonnegativity bounds.


\vspace{5mm}

Now suppose instead that the logistic loss is used:
\[
\ell(y_i,\beta^\top x_i+\beta_0)
=
\log\left(
1+e^{-y_i(\beta^\top x_i+\beta_0)}
\right).
\]

Then (2.2) contains constraints of the form
\[
\theta+u_i
\geq
\log\left(
1+e^{-y_i(\beta^\top x_i+\beta_0)}
\right).
\]

Unlike hinge loss, logistic loss is a \textbf{smooth nonlinear convex function},
rather than a piecewise-linear function. Therefore, it cannot be
represented exactly as a linear mixed-integer optimization problem using
only finitely many binary variables and linear constraints.

However, binary variables are not necessary. As in Lecture 6 (and in Problem 1 part (e)), applying
the relaxation and strong-duality argument leads to a continuous convex
nonlinear optimization problem. Thus, although the formulation is not
an LP or an exact linear MIO, it is a convex nonlinear optimization
problem and is computationally tractable.


\vspace{5mm}

\subsection*{(d)}


\textbf{(i)}

At a solution $(\beta^t,\beta_0^t,t^t)$, we define the current
losses as
\[
c_i^t
=
\ell\left(y_i,(\beta^t)^\top x_i+\beta_0^t\right),
\qquad i=1,\ldots,n.
\]

The separation problem is therefore 
\[
\max_{z\in\mathcal Z}
\sum_{i=1}^n z_i c_i^t.
\]

By part (a), this problem has a closed-form solution: the optimal $z$
selects the $k$ observations with the largest losses. If
\[
c_{(1)}^t\geq c_{(2)}^t\geq\cdots\geq c_{(n)}^t,
\]
then the optimal value of the separation problem is
\[
\sum_{j=1}^k c_{(j)}^t.
\]

So if we have
\[
\sum_{j=1}^k c_{(j)}^t>t^t,
\]
our corresponding $z^*$ gives a violated constraint, so we add the cut
\[
t
\geq
\sum_{i=1}^n
z_i^*
\ell(y_i,\beta^\top x_i+\beta_0)
\]
and re-solve the restricted problem. Otherwise, no constraint is
violated and the current solution is optimal.

Computing the $n$ losses requires $O(np)$ time, since evaluating each
$\beta^\top x_i$ requires $O(p)$ operations. Sorting the losses to find
the $k$ largest requires $O(n\log n)$ time. Thus, a separation step can
be performed in
\[
\boxed{O(np+n\log n)}
\]
time.

\textbf{(ii)}

One reason to prefer the cutting-plane approach is scalability.
Formulation (2.2) introduces $n$ additional variables $u_i$ and at least
$n$ additional constraints, so its size grows directly with the number
of observations.

However, our cutting-plane master problem has only
\[
p+2
\]
variables, which are the $p$ components of $\beta$, $\beta_0$, and $t$. It starts
with only a small subset of the constraints in (2.3) and adds a new
constraint only when the separation problem finds a violated one.
Therefore, if relatively few cuts are needed, the optimization problems
solved at each iteration can be much smaller than solving (2.2)
directly.


\vspace{8mm}

\section*{Question 3}


\subsection*{Part (a)}

\textbf{(i)}

Imposing the four candidate constraint sets on every adjacent pair gave
the following results:
\[
\begin{array}{c|c|c|c}
\text{Candidate}
& \text{Constraints imposed}
& \text{Objective}
& \text{Broken constraints}\\
\hline
\text{Forward} & 39 & 1.0528 & 574\\
\text{Backward} & 39 & 1.4079 & 506\\
\text{Both} & 78 & 2.4454 & 0\\
\text{Stepped backwards} & 39 & 1.0528 & 908
\end{array}
\]

Thus, imposing \textbf{both the forward and backward constraints} on every
adjacent pair is the only candidate that leaves none of the full 1560
convexity constraints violated. Since there are $n-1=39$ adjacent pairs
and two constraints are imposed on each pair, this formulation uses
$
2(n-1)=78
$
constraints.

Its objective value is
$
2.4454,
$
which is exactly the same as the objective obtained using all 1560
constraints. Thus, for this one-dimensional problem, the 78 adjacent-pair constraints are sufficient. 


\vspace{5mm}
\textbf{(ii)}

For the random control sets, increasing the number of constraints
generally moved the solution closer to the fully convex fit. The results
were
\[
\begin{array}{c|c|c}
\text{Random constraints}
& \text{Objective}
& \text{Broken constraints}\\
\hline
39  & 0.5533 & 608\\
78  & 0.7590 & 517\\
156 & 1.3073 & 380\\
312 & 1.6723 & 248\\
624 & 2.2773 & 70
\end{array}
\]

As more random constraints are added, the objective gets closer to the
full-problem objective of $2.4454$, and fewer constraints are violated.
However, having more constraints does not necessarily mean they are the
right ones. For example, 78 random constraints still leave 517
violations, while the 78 bidirectional adjacent-pair constraints from
part (i) leave no violations.

This is because the adjacent-pair constraints directly enforce the
ordering of slopes needed for convexity. Random constraints can miss
these important neighboring relationships. So we should choose the right constraints instead of choosing more of them. 

\newpage


\textbf{(iii)}

From part (i), we impose both the forward and backward constraints on
every adjacent pair $(k,k+1)$.

\textbf{(1)}

We impose on a single adjacent pair $(k, k+1)$. The forward constraint is
\[
\theta_k+\xi_k(x_{k+1}-x_k)\leq\theta_{k+1}.
\]
Since
\[
h_k=x_{k+1}-x_k>0,
\]
we can rewrite this as
\[
\xi_k h_k\leq\theta_{k+1}-\theta_k.
\]

\[
\xi_k
\leq
\frac{\theta_{k+1}-\theta_k}{h_k}
=
c_k.
\]

The backward constraint is
\[
\theta_{k+1}+\xi_{k+1}(x_k-x_{k+1})\leq\theta_k.
\]
Since $x_k-x_{k+1}=-h_k$, this becomes
\[
\theta_{k+1}-\xi_{k+1}h_k\leq\theta_k.
\]
We rearrange and get: 
\[
c_k\leq\xi_{k+1}.
\]

\[
\xi_k\leq c_k\leq\xi_{k+1}.
\]

Neither inequality alone is sufficient. The forward constraint only gives
$\boxed{\xi_k\leq c_k}$, while the backward constraint only gives
$\boxed{c_k\leq\xi_{k+1}}$. Both are needed to connect the slopes on the two
sides.


\vspace{5mm}
\textbf{(2)}

Now we apply this to every adjacent pair.

The result above holds for every $k=1,\ldots,n-1$. Therefore, the
inequalities can be placed end to end:
\[
\xi_1
\leq c_1
\leq \xi_2
\leq c_2
\leq \cdots
\leq \xi_{n-1}
\leq c_{n-1}
\leq \xi_n.
\]

Here, we see how both the subgradients and the chord slopes are nondecreasing, i.e.:
\[
\boxed{\xi_1\leq\xi_2\leq\cdots\leq\xi_n,}
\]
and
\[
\boxed{c_1\leq c_2\leq\cdots\leq c_{n-1}.}
\]

This is where the ordering $x_1<x_2<\cdots<x_n$ is important. Since
$x_{k+1}>x_k$, we have $h_k=x_{k+1}-x_k>0$, so dividing by $h_k$
does not reverse the direction of the inequalities. Also,  we see how the
adjacent pairs are $(1,2),(2,3),\ldots,(n-1,n)$, so these consecutive
inequalities share a common subgradient and can be linked together in one ordered sequence (like we see above). 


\vspace{5mm}
\textbf{(3)}

Now we consider two nonadjacent points.

First suppose $i<j$. We want to show that the original convexity
constraint
\[
\theta_i+\xi_i(x_j-x_i)\leq\theta_j
\]
holds even though it was not explicitly imposed.

We are given the height difference and distance telescope are
\[
\theta_j-\theta_i
=
\sum_{k=i}^{j-1}c_kh_k,
\]

\[
x_j-x_i
=
\sum_{k=i}^{j-1}h_k.
\]

From the ordering in part (2), for every $k=i,\ldots,j-1$,
\[
\xi_i\leq c_k.
\]
Since $h_k>0$,
\[
\xi_i h_k\leq c_kh_k.
\]
We sum these inequalities: 
\[
\xi_i\sum_{k=i}^{j-1}h_k
\leq
\sum_{k=i}^{j-1}c_kh_k.
\]

Now we use the telescoping expressions: 
\[
\xi_i(x_j-x_i)
\leq
\theta_j-\theta_i.
\]
Therefore,
\[
\theta_i+\xi_i(x_j-x_i)\leq\theta_j.
\]

Now suppose $i>j$. We must also prove the reverse-direction constraint
\[
\theta_i+\xi_i(x_j-x_i)\leq\theta_j.
\]

As before, we use the telescoping expressions,
\[
\theta_i-\theta_j
=
\sum_{k=j}^{i-1}c_kh_k,
\qquad
x_i-x_j
=
\sum_{k=j}^{i-1}h_k.
\]

The difference from the $i<j$ case is that the ordering from part (2)
now gives us 
\[
c_k\leq\xi_i,
\qquad k=j,\ldots,i-1.
\]
Since $h_k>0$, when we sum over these intervals we get 
\[
\theta_i-\theta_j
=
\sum_{k=j}^{i-1}c_kh_k
\leq
\xi_i\sum_{k=j}^{i-1}h_k
=
\xi_i(x_i-x_j).
\]
We just rearrange and get the result: 
\[
\theta_i+\xi_i(x_j-x_i)\leq\theta_j.
\]

Thus, the $i<j$ case proves the constraints from a point to every point
to its right, while the $i>j$ case proves the constraints to every point
to its left. Since the two cases cover both directions between any pair of points,
we have shown that every constraint in the full formulation is
automatically satisfied. Therefore, we only need the $2(n-1)$
bidirectional adjacent-pair constraints instead of explicitly imposing
all $n(n-1)$ constraints. \qed


\newpage
\subsection*{Part (b)}

\textbf{(i)}

The worst violation is caused by the supporting plane at point $28$,
which passes above point $12$. The size of the violation is approximately 0.518.


\vspace{5mm}
\textbf{(ii)}

The step that fails is \textbf{Step (2)} of the proof in part (a)(iii). In one
dimension, the points can be ordered as
\[
x_1<x_2<\cdots<x_n,
\]
and the adjacent-pair constraints give
\[
\xi_k\leq c_k\leq\xi_{k+1}.
\]
These inequalities can then be linked together into the single chain
\[
\xi_1\leq c_1\leq\xi_2\leq c_2\leq\cdots\leq\xi_n.
\]

In $\mathbb{R}^2$, however, there is no single left-to-right ordering of
the points. The Delaunay neighbours form a two-dimensional mesh rather
than a chain. 

Also, in one dimension each $\xi_i$ is just one number, so it makes
sense to compare two slopes and write something like
\[
\xi_i\leq\xi_j.
\]
In two dimensions, however, each $\xi_i$ has two components, so there is no
single left-to-right slope ordering that works in the same way.

Because of this, the inequalities from neighboring points cannot be
linked together like they were in part (a). As a result, checking all
Delaunay neighbors in both directions is not enough to guarantee the
constraints between points that are farther apart.


\vspace{5mm}

\subsection*{Part (c)}


\textbf{(i)}

Starting with both directions, it took \textbf{$2$ rounds} to reach no broken
constraints. The final model had \textbf{$231$ constraints} and an objective
value of 2.58541.

Starting with only one direction per edge, it took \textbf{$6$ rounds} to reach
no broken constraints. The final model had \textbf{$207$ constraints} and the
same objective value of 2.58541.


So, starting with both directions took fewer rounds since we started
with more constraints already included. Starting with one direction
took longer, but it ended up using fewer constraints overall. In the
end, both starting points reached the same objective value.


\vspace{5mm}
\textbf{(ii)}

If $n=10{,}000$, the full formulation would contain
\[
n(n-1)
=
10{,}000(9{,}999)
=
99{,}990{,}000
\]
constraints.

This is almost $100$ million constraints, so solving the full model
directly would be much more difficult. The purpose of constraint
generation is to avoid including all of these constraints from the
start. Instead, we begin with a smaller set and only add constraints
that are actually violated.




\end{document}
